% ---------------------------------------------------------------------------
%  Stile comune di esercitazione e soluzioni (pdfLaTeX).
%  Viene incluso da esercitazione.tex e soluzioni.tex con \input.
% ---------------------------------------------------------------------------
\documentclass[11pt,a4paper]{article}
\usepackage[a4paper, top=14mm, bottom=18mm, left=16mm, right=16mm, footskip=9mm]{geometry}
\usepackage[utf8]{inputenc}
\usepackage[T1]{fontenc}
\usepackage{helvet}
\renewcommand{\familydefault}{\sfdefault}
\usepackage{amsmath, amssymb, graphicx, xcolor, fancyhdr, lastpage, tabularx, booktabs, array, enumitem}
\usepackage[most]{tcolorbox}
\usepackage[siunitx, american resistors, american currents]{circuitikz}
\ctikzset{bipoles/thickness=1, current/distance=0.35}
\usetikzlibrary{calc, arrows.meta}

\definecolor{titolo}{HTML}{2A7FA8}
\definecolor{grigio}{HTML}{E9ECEF}
\definecolor{tenue}{HTML}{555555}
\setlength{\parindent}{0pt}
\setlength{\parskip}{3pt}

\newcommand{\U}[1]{\,\mathrm{#1}}
\newcommand{\virg}{{,}}
\newcommand{\riga}[1][28mm]{\underline{\hspace{#1}}}

% intestazione come nelle esercitazioni della scuola
\newcommand{\intestazione}[2]{%
  \begin{minipage}[c]{20mm}\includegraphics[width=18mm]{logo-isis.png}\end{minipage}%
  \begin{minipage}[c]{\dimexpr\linewidth-20mm}
    {\fontsize{20}{24}\selectfont\bfseries\color{titolo} #1}\par\vspace{1mm}
    {\footnotesize\bfseries Indirizzo: IP14 -- MANUTENZIONE E ASSISTENZA TECNICA \quad -- \quad Classe 3MAT}
  \end{minipage}\par\vspace{2mm}
  {\color{titolo}\rule{\linewidth}{0.8pt}}\par\vspace{3mm}
  \begin{tcolorbox}[colback=grigio, colframe=grigio, arc=0pt, left=6pt, top=4pt, bottom=4pt]
    {\Large\bfseries #2}
  \end{tcolorbox}}

\newcommand{\nota}[1]{\begin{tcolorbox}[colback=grigio!45, colframe=grigio, arc=0pt, left=6pt, top=3pt, bottom=3pt]\small #1\end{tcolorbox}}
\newcommand{\tipo}[1]{\par\vspace{2mm}{\large\color{titolo} #1}\par\vspace{-1mm}{\color{titolo!60}\rule{\linewidth}{0.4pt}}\par\vspace{1mm}}
\newcommand{\esercizio}[2]{\par\vspace{3mm}\textbf{Esercizio #1.} \textbf{#2}\par}

% alimentazione: barra con etichetta
\newcommand{\rail}[2]{\draw[line width=1.2pt] ($#1+(-0.6,0)$) -- ($#1+(0.6,0)$); \node[above] at #1 {#2};}

\pagestyle{fancy}
\fancyhf{}
\renewcommand{\headrulewidth}{0pt}
\newcommand{\piede}[1]{\fancyfoot[L]{\scriptsize I.S.I.S. Sacile-Brugnera | sede: ``Piero Della Valentina'' | #1}
  \fancyfoot[R]{\scriptsize Pagina \thepage\ di \pageref{LastPage}}}
